\documentclass[12pt,a4paper]{article}
\usepackage[utf8]{inputenc}
\usepackage[portuguese]{babel}
\usepackage[margin=2.5cm]{geometry}
\usepackage{hyperref}
\usepackage{enumitem}
\usepackage{titlesec}
\usepackage{xcolor}

% Configuração de cores e links (Estilo Canadá)
\definecolor{canada-red}{RGB}{191, 33, 47}
\hypersetup{
    colorlinks=true,
    linkcolor=canada-red,
    urlcolor=canada-red
}

\titleformat{\section}{\normalfont\Large\bfseries\color{canada-red}}{\thesection}{1em}{}
\titleformat{\subsection}{\normalfont\large\bfseries}{\thesubsection}{1em}{}

\title{\textbf{Pós-Graduação com Bolsa no Canadá}}
\author{}
\date{}

\begin{document}

\maketitle

\textbf{Sim, existe!} O Canadá possui um sistema de financiamento muito robusto para pós-graduações, especialmente para programas \textit{stricto sensu} (Mestrado e Doutorado). 

No entanto, a lógica de bolsas no Canadá funciona de maneira diferente de outros países. Em vez de uma ``bolsa única'' centralizada, o financiamento canadense costuma ser composto por um \textbf{pacote de apoio financeiro (Funding Package)}, que combina prêmios por mérito, bolsas governamentais e trabalho dentro da universidade.

Abaixo, veja como se dividem as principais oportunidades para estudantes internacionais:

\section{Bolsas Internas das Universidades (Graduate Funding Packages)}
Ao ser aceito em um \textbf{Mestrado Acadêmico (Thesis-based Master's)} ou \textbf{Doutorado (PhD)} em uma universidade pública canadense, é muito comum que você receba uma oferta de financiamento junto com a sua carta de aceitação. Esse pacote geralmente exige que você trabalhe no campus e inclui:

\begin{itemize}[leftmargin=*]
    \item \textbf{Teaching Assistantship (TA):} Você atua como assistente de ensino, ajudando professores a corrigir provas, aplicar trabalhos ou ministrar turmas de graduação. É uma atividade remunerada (com salários competitivos).
    \item \textbf{Research Assistantship (RA):} Você é pago diretamente pelo orientador (através de verbas de pesquisa dele) para trabalhar em um projeto de laboratório ou coletar dados que geralmente têm relação direta com a sua própria tese.
    \item \textbf{Graduate Fellowships / Entrance Scholarships:} Prêmios em dinheiro concedidos diretamente pela faculdade com base exclusivamente no seu mérito acadêmico (histórico, currículo e cartas de recomendação), sem necessidade de contrapartida de trabalho.
\end{itemize}

\begin{quote}
    \textbf{Atenção:} Mestrados profissionais (voltados apenas para o mercado, os chamados \textit{Course-based} ou \textit{Non-thesis}, como a maioria dos MBAs) raramente oferecem esses pacotes de financiamento. Eles costumam oferecer apenas pequenas bolsas de entrada (\textit{Entrance Awards}) de 2.000 a 5.000 CAD.
\end{quote}

\section{Bolsas de Prestígio do Governo Canadense}
O governo federal e os governos provinciais do Canadá oferecem bolsas de altíssimo valor para atrair talentos mundiais. As mais famosas e generosas são:

\begin{itemize}[leftmargin=*]
    \item \textbf{Vanier Canada Graduate Scholarships:} Focada exclusivamente em \textbf{Doutorado (PhD)}. É uma das bolsas mais concorridas do mundo, pagando \textbf{50.000 CAD por ano}, por até 3 anos. Avalia excelência acadêmica, potencial de pesquisa e liderança.
    \item \textbf{Canada Graduate Research Scholarships - Doctoral (CGRS-D):} Financiamento de peso administrado pelas três agências federais de pesquisa do Canadá (nas áreas de saúde, engenharias/ciências naturais e ciências humanas), pagando em média \textbf{40.000 CAD por ano}.
    \item \textbf{Bolsas Provinciais:} Governos de províncias específicas oferecem incentivos próprios, como a \textit{Ontario Graduate Scholarship (OGS)} em Ontário e a \textit{Alberta Graduate Excellence Scholarship (AGES)} em Alberta. Elas giram em torno de 15.000 CAD anuais e muitas são abertas a estrangeiros.
\end{itemize}

\section{Programas de Cooperação Internacional (Para Brasileiros)}
Existem acordos específicos que financiam estudantes da América Latina:

\begin{itemize}[leftmargin=*]
    \item \textbf{ELAP (Emerging Leaders in the Americas Program):} Ideal para quem já está matriculado em uma pós-graduação no Brasil e quer fazer um \textbf{intercâmbio de pesquisa (doutorado ou mestrado sanduíche)} de 4 a 6 meses no Canadá. A bolsa cobre passagens, visto e custos de vida enquanto você estiver lá.
    \item \textbf{Mitacs Accelerate International:} Oferece subsídios para estágios de pesquisa aplicada de curta e média duração em universidades canadenses em parceria com indústrias e empresas locais.
\end{itemize}

\section{Como aumentar suas chances de conseguir?}

\begin{enumerate}
    \item \textbf{Encontre um Orientador Primeiro:} Para programas acadêmicos e de pesquisa, a aprovação da bolsa muitas vezes depende de um professor aceitar ser seu supervisor. Entre em contato com professores cujas linhas de pesquisa alinhem-se perfeitamente com o seu histórico acadêmico ou projetos anteriores.
    \item \textbf{Capriche no Projeto e no Histórico:} As bolsas canadenses são majoritariamente baseadas em \textbf{mérito}. Notas altas na graduação, publicações acadêmicas e um projeto de pesquisa inovador são fundamentais.
    \item \textbf{Prepare a Documentação com Antecedência:} Você precisará de exames de proficiência em inglês (IELTS/TOEFL) ou francês, cartas de recomendação acadêmicas fortes e tradução juramentada dos seus históricos escolares.
\end{enumerate}

Para pesquisar opções específicas abertas no momento atual para o seu perfil e país de origem, o portal oficial \textbf{EduCanada} (\href{https://www.educanada.ca}{educanada.ca}) possui um buscador excelente de bolsas internacionais estruturado pelo próprio governo canadense.

\end{document}